\documentclass[11pt, letterpaper]{article}
% Packages
\usepackage[
    ignoreheadfoot,
    top=1.5cm, bottom=1.5cm, left=1.8cm, right=1.8cm,
    footskip=0.8cm
]{geometry}
\usepackage{enumitem}
\usepackage{hyperref}
\usepackage{titlesec}
\usepackage{tabularx}
\usepackage{needspace}
\usepackage[usenames,dvipsnames]{xcolor}
\usepackage{fontawesome5}
\hypersetup{hidelinks}
% Colors & Fonts
\definecolor{primaryColor}{RGB}{0, 0, 0} % Black for objectivity
\urlstyle{same}
% Formatting
\pagestyle{empty}
\raggedbottom
\setlength{\parindent}{0pt}
\setlength{\tabcolsep}{0in}
% Section formatting
\titleformat{\section}{
  \vspace{-4pt}\scshape\raggedright\large
}{}{0em}{}[\color{black}\titlerule \vspace{-5pt}]
% Custom commands
\newcommand{\resumeItem}[1]{
  \item\small{
    {#1 \vspace{-2pt}}
  }
}
\newcommand{\resumeSubheading}[4]{
  \vspace{-1pt}\item
    \begin{tabular*}{\linewidth}[t]{@{\extracolsep{\fill}}lr@{}}
      \textbf{#1} & #2 \\
      \textit{\small#3} & \textit{\small #4} \\
    \end{tabular*}\vspace{-5pt}
}
\newcommand{\resumeProjectHeading}[2]{
    \Needspace{4\baselineskip}
    \item
    \begin{tabularx}{\linewidth}{@{}>{\raggedright\arraybackslash}Xr@{}}
      \small#1 & #2 \\
    \end{tabularx}\vspace{-5pt}
}
\newcommand{\resumeSubItem}[1]{\resumeItem{#1}\vspace{-4pt}}
\renewcommand\labelitemi{$\vcenter{\hbox{\tiny$\bullet$}}$}
\begin{document}
%----------HEADER----------
\begin{center}
    \textbf{\Huge \scshape Guandi Wang} \\ \vspace{4pt}
    \small
    \faIcon{map-marker-alt} Stockholm, Sweden \quad
    \faIcon{envelope} \href{mailto:guandi@kth.se}{guandi@kth.se} \quad
    \faIcon{phone} +46 760219132 \quad
    \faIcon{linkedin} \href{https://www.linkedin.com/in/guandi-wang-w}{linkedin.com/in/guandi-wang-w}
\end{center}
%-----------EDUCATION-----------
\section{Education}
\begin{itemize}
  \resumeSubheading
    {EURECOM}{Sophia Antipolis, France}
    {M.Sc. in Sensing Big Data for Intelligent Robots}{July 2026 -- Present}
  \resumeSubheading
    {KTH Royal Institute of Technology}{Stockholm, Sweden}
    {M.Sc. in ICT Innovation: Autonomous Systems and Intelligent Robots}{Aug. 2025 -- Present}
    \resumeItem{Selected Coursework: Machine Learning, Distributed Artificial Intelligence, Introduction to Robotics, Hybrid and Embedded Control Systems.}
  \resumeSubheading
    {Xi'an University of Architecture and Technology}{Shaanxi, China}
    {B.Sc. in Building Electrical and Intelligence}{Sep. 2021 -- July 2025}
    \resumeItem{GPA: 83.5/100 (Top 10\%) $|$ \textbf{Technical Focus:} Control Theory, Power Engineering, Building Environment, Embedded Systems.}
\end{itemize}
%-----------RESEARCH EXPERIENCE-----------
\section{Research Experience}
  \begin{itemize}[leftmargin=0.15in, label={}]
    % --- Tekniska Mekanik, KTH ---
    \resumeSubheading
      {Research Intern}{Sep. 2025 -- Present}
      {Tekniska Mekanik, KTH Royal Institute of Technology}{Supervisor: Junle Liu}
    \resumeProjectHeading
      {\textbf{ForeWind: Physics-Informed Spatiotemporal Wind Flow Forecasting}}{}
      \resumeItem{Constructed a 26-year hourly ERA5 benchmark (1999--2024) for wind forecasting over heterogeneous terrain. Built a physics-informed feature space from mean wind flow, wind variability, diurnal temperature range, elevation, and terrain roughness; used a Self-Organizing Map to identify nine wind regimes and select 12 representative coastal, inland, and plateau cities.}
      \resumeItem{Developed a unified data and training pipeline for eight deep-learning architectures on city-centered $64 \times 64$ patches under a 24-hour input-to-24-hour forecast protocol, yielding 96 city--model configurations on Slurm-based HPC systems. Evaluated city RMSE, temporal correlation, patch-scale anomaly correlation, and persistence-normalized skill.}
      \resumeItem{Showed that lower-boundary setting governs attainable error scale while architecture controls regime-specific failure modes. Two-way fixed-effect decomposition attributed most RMSE variability to city setting and persistence-normalized skill to city--model interaction; grid-level skill maps, lead-dependent errors, and power spectral density analysis exposed coherent variance attenuation across coastal and plateau regimes. \href{https://github.com/LIUJUNLE97/forewind}{Code and data}.}

    \resumeProjectHeading
      {\textbf{Physics-Constrained Generative Reconstruction of Incompressible Flows}}{}
      \resumeItem{Investigating how generative flow-reconstruction models can strictly satisfy incompressibility through divergence-free constraints while preserving high reconstruction accuracy and calibrated uncertainty quantification.}
      \resumeItem{Combining Gaussian Splatting with diffusion models to generate physically consistent flow fields on datasets including the Johns Hopkins Turbulence Database (JHTDB) and Taylor--Green vortex (TGV) benchmarks.}

    \resumeProjectHeading
      {\textbf{CFD-Informed Reinforcement Learning Environments for UAVs}}{}
      \resumeItem{Building two Isaac Lab UAV simulation environments from CFD data: a time-varying wind-field scenario based on JHTDB and a built-environment scenario driven by CFD flow around a Y-shaped building.}
      \resumeItem{Designed the simulation lab for HPC execution and integrated ParaView with ROS2 to support real-time visualization and video streaming of large-scale aerodynamic environments.}

    \resumeProjectHeading
      {\textbf{GPU-Accelerated OpenFOAM Deployment on HPC}}{}
      \resumeItem{Deploying and validating GPU-enabled OpenFOAM workflows on supercomputing infrastructure for scalable CFD simulation and downstream learning tasks.}

  \end{itemize}
\newpage
\section{Research Experience (Continued)}
  \begin{itemize}[leftmargin=0.15in, label={}]
    \resumeProjectHeading
      {\textbf{Deep Multimodal Object Detection} $|$ \emph{Collaboration with Junle Liu}}{}
      \resumeItem{Developed spatial-mask interaction and channel-competition mechanisms for robust RGB--thermal object detection under sensor misalignment and adverse conditions, with an emphasis on deployable multimodal perception for real robotic systems. Accepted to BMVC 2026.}

    % --- IIT Guest Student ---
    \resumeSubheading
      {Humanoid Sensing and Perception Guest Student}{June 2026 -- Sep. 2026}
      {Istituto Italiano di Tecnologia (IIT)}{Supervisor: Maria Lombardi; PI: Lorenzo Natale}
    \resumeProjectHeading
      {\textbf{Unified Visual Segmentation and Annotation Tool}}{}
      \resumeItem{Developed a visual segmentation and annotation tool using SAM3.1 as the foundation model and integrating RTMPose and RAFT for pose-aware annotation and optical-flow propagation.}
      \resumeItem{Supported precise frame-by-frame annotation, automatic object tracking across video sequences, and manual correction for both live-action movies and animated content.}

    \resumeProjectHeading
      {\textbf{Force-Aware Physical Interaction with the R1 Humanoid Robot}}{}
      \resumeItem{Developed a teleoperation-based physical interaction setup in which a human and the R1 humanoid robot jointly grasp an interaction-force sensor, enabling synchronized motion and force exchange for studying human--robot mechanical interaction.}

  \end{itemize}
%-----------SELECTED PROJECTS-----------
\section{Selected Projects}
  \begin{itemize}[leftmargin=0.15in, label={}]
    % --- Project 3: ROS2 Navigation (KTH Coursework) ---
    \resumeProjectHeading
      {\textbf{Autonomous Mobile Robot Navigation System} (Coursework) }{Sep. 2025 -- Dec. 2025}
      \resumeItem{Architected a navigation stack for TurtleBot3 using \textbf{ROS2 (Foxy)} and Python. Migrated control logic from Finite State Machines to \textbf{Behavior Trees} for enhanced reactivity.}
     
      \resumeItem{Integrated AMCL for robust localization and tuned PID controllers for smooth trajectory tracking.}
    % --- Project 4: Bachelor Thesis ---
    \resumeProjectHeading
      {\textbf{Photovoltaic Panel Defect Detection Algorithm} $|$ \emph{Bachelor Thesis}}{Jan. 2025 -- Jun. 2025}
      \resumeItem{\textbf{Grand Prize}, Undergraduate Graduation Project Competition, Shaanxi Automation Association \hfill 2025.
     }
    
  \end{itemize}
%-----------PUBLICATIONS-----------
\section{Publications}
  \begin{itemize}[leftmargin=0.15in, label={}]
    \resumeItem{\textbf{Deep Multimodal Object Detection via Spatial Mask Interaction and Channel Competition.} Accepted to BMVC 2026. \href{https://arxiv.org/abs/2608.02092}{arXiv:2608.02092}.}
  \end{itemize}
%-----------SKILLS-----------
\section{Technical Skills}
 \begin{itemize}[leftmargin=0.15in, label={}]
    \small{\item{
     \textbf{Languages}{: Python, C/C++, MATLAB, LaTeX} \\
     \textbf{Frameworks \& Libraries}{: PyTorch, ROS2 (Foxy/Humble), OpenCV, Scikit-learn, Isaac Lab, YARP, LeRobot} \\
     \textbf{Tools \& Platforms}{: Linux (Ubuntu), Git, Docker, Gazebo, Slurm (HPC), OpenFOAM, ParaView, VR systems}
  
    }}
 \end{itemize}
%-----------AWARDS-----------
\section{Selected Awards}
 \begin{itemize}[leftmargin=0.15in, label={}]
    \small{
     \item{ \textbf{Outstanding Graduate}, Xi'an University of Architecture and Technology \hfill 2025}
    \item{
     \textbf{Academic Scholarship}, Class II, Awarded for 3 consecutive years \hfill 2021 -- 2023
    }}
   
    
 \end{itemize}
\end{document}
